% TOP_PROBLEM: {"id": 2302043, "title": "Research Problems in Function Theory — Problem 2.43", "category_id": 8, "category": {"id": 8, "name": "analysis", "display_name": "Analysis", "description": "Limits, continuity, calculus, and function theory.", "slug": "analysis", "order_index": 8, "created_at": "2026-07-31T15:26:25.671Z"}, "status": "open"}
% FIRST_POSTED: null
% ENTRY_KIND: statement_only
% Imported from: batch11.tex; UnsolvedMath ID 2302043
% Compile twice: lualatex 2302043.tex
\documentclass[11pt]{article}
\usepackage{amsmath,amssymb,mathtools}
\usepackage{fontspec}
\usepackage{unicode-math}
\setmainfont{Latin Modern Roman}
\setmathfont{Latin Modern Math}
\usepackage{xurl}
\usepackage[unicode,colorlinks=true,linkcolor=blue,urlcolor=blue,pdftitle={UnsolvedMath Problem 2302043}]{hyperref}
\newcommand{\sref}[2]{\hyperlink{source:#1:#2}{[S#2]}}
\newcommand{\cataloguescope}[1]{\paragraph{Mathematical statement.}#1}
\begin{document}
% UnsolvedMath ID 2302043; source catalogue code AMR-022-2043
\setcounter{section}{2302042}
\section{Research Problems in Function Theory — Problem 2.43}\label{2302043}
{\small\sffamily UnsolvedMath ID 2302043 · Analysis}
\subsection{Definitions and mathematical statement}

\paragraph{Shared notation.}$\mathbb N=\{1,2,\ldots\}$, and $\mathbb Z,\mathbb Q,\mathbb R,\mathbb C$ denote the integers, rationals, reals and complex numbers. Any other conventions are specified locally.


For an entire function $f$, write
\[
 M(r,f)=\max_{|z|=r}|f(z)|,\qquad m_0(r,f)=\min_{|z|=r}|f(z)|.
\]
For nonconstant $f$, its order is $\rho(f)=\limsup_{r\to\infty}\log\log M(r,f)/\log r$;
the lower order is defined using $\liminf$ instead. Every constant entire function, including the zero function, has order and lower order zero. In logarithmic modulus ratios, use the extended-real convention $\log0=-\infty$.
A function permutes a set $S$ if $f|_S:S\to S$ is bijective. For an entire function of order one, its exponential type is $\tau(f)=\limsup_{r\to\infty}\log M(r,f)/r$.

\paragraph{Statement recorded in the supplied source.}
If $f$ is transcendental entire and permutes $\mathbb Z$, must $\rho(f)\ge1$, with $\tau(f)\ge\pi$ whenever $\rho(f)=1$? Ask the same question when $f$ permutes $\mathbb N=\{1,2,\ldots\}$. \sref{2302043}{2}
Update 2.43 supplies the counterexample $f(z)=z+4\cos(\pi(z-2)/3)$, which permutes both sets and has order one and type $\pi/3$. It also asks for examples of reasonably small type that substantially rearrange the positive integers, without specifying a quantitative scrambling criterion. \sref{2302043}{2}
\subsection{Short English statement}
Do entire functions that bijectively rearrange either all integers or the positive integers require the proposed order and exponential type?
\subsection{Sources}
\begin{description}
\item[\hypertarget{source:2302043:1}{S1}] ULAM.AI and the UnsolvedMath Contributors, \emph{UnsolvedMath: A Curated Collection of Open Mathematics Problems}, record 2302043 (AMR-022-2043). CC BY 4.0. \url{https://huggingface.co/datasets/ulamai/UnsolvedMath}.
\item[\hypertarget{source:2302043:2}{S2}] W. K. Hayman and E. F. Lingham, \emph{Research Problems in Function Theory} (2018), Chapter 2, Problem 2.43 and Update 2.43. \url{https://arxiv.org/abs/1809.07200}.
\end{description}
\label{end:2302043}
\end{document}
